% !TeX root = ../main.tex
% Ground movement, monitoring and the observational method. Merged:
% predicting movement, assessing third-party impact, setting limits and
% monitoring against them are one workflow, not two subjects.
\chapter{Ground Movement, Monitoring and the Observational Method}
\label{ch:movement-monitoring}
\label{ch:settlement-damage}
\label{ch:monitoring}

In urban tunnelling the governing constraints are commonly imposed by assets
the designer does not control. Buildings, railways, utilities and existing
tunnels carry movement limits far below any threshold that would endanger the
excavation, and satisfying them determines the excavation method, the face
pressure, the extent of ground treatment and, on occasion, the alignment.
Prediction and monitoring are treated here as one subject. The chain runs from
construction control through ground movement and structural strain to damage;
every link carries uncertainty; and the only defence against that uncertainty
is to predict, set limits, measure, and act on the measurement. That sequence
is the observational method, and third-party protection is its most demanding
application.

\section{Predicting ground movement}
\label{sec:predicting-movement}

\subsection{The ground movement field}
\label{sec:movement-field}

Surface settlement above a tunnel is described by a Gaussian curve whose
maximum $S_{\max}$ lies above the axis and whose width $i$, the distance from
the centreline to the point of inflection, is set by the trough width
parameter $K$ and the depth to the axis $z_0$ through $i = K z_0$:
\[
S_y = S_{\max} \exp\!\left(-\frac{y^2}{2i^2}\right),
\]
with $S_y$ the settlement at transverse distance $y$ from the centreline.
Section~\ref{sec:gaussian-trough} derives $S_{\max}$ from the volume loss and
gives the customary values of $K$, and Section~\ref{sec:seq-settlement} sets
out the calculation sequence. The area under the curve equals the volume loss ---
the ground that entered the tunnel and did not come out as intended
excavation, expressed as a proportion of the theoretical excavated volume.
Volume loss is a construction outcome, not a ground property (Table~\ref{tab:volume-loss}): the same ground
gives very different volume loss under different face pressure and grouting
control.

The transverse Gaussian is only one section through a three-dimensional field.
Longitudinally, settlement develops as the face approaches and passes,
following a cumulative probability form: some settlement occurs ahead of the
face, roughly half is complete as the face passes beneath, and the remainder
develops behind. A structure therefore experiences a settlement history, not a
single value. It is subjected to hogging as the face approaches, arching over
the advancing edge of the trough, and to sagging after it passes, as the
trough deepens beneath it --- a reversal that a static assessment does not
capture.

Subsurface movements matter for anything founded below the surface. The trough
narrows and deepens with depth towards the tunnel, so a piled foundation may
experience considerably more movement than the surface above it. In clays the
trough width at depth $z$ below the surface is approximately
\[
i = 0.175\, z_0 + 0.325\,(z_0 - z),
\]
which recovers $i = 0.5\,z_0$ at the surface~\citep{mair1993subsurface}. Taking
the trough volume $V_s$ as unchanged with depth, as for undrained clay, the
narrower trough gives a proportionally larger maximum,
$S_{\max} = V_s/(i\sqrt{2\pi})$. Horizontal
movement accompanies the vertical: if the displacement vectors are taken to
point towards the tunnel axis, the horizontal displacement at transverse
distance $y$ is $S_h = (y/z_0)\, S_y$, and differentiating it gives the
horizontal strain~\citep{mair1996subsurface}
\[
\varepsilon_h = \frac{S_y}{z_0}\left(\frac{y^2}{i^2} - 1\right),
\]
taken positive in tension, with $S_y$ the settlement at $y$ from the
Gaussian expression. The strain is compressive within the points of
inflection, reaching $S_{\max}/z_0$ over the centreline, and tensile beyond
them, reaching $2e^{-3/2}S_{\max}/z_0 \approx 0.45\,S_{\max}/z_0$ at
$y = \sqrt{3}\,i$. The tensile zone at the trough edge is where most building
damage occurs, and an assessment that considers settlement alone will miss it.

\subsection{Building damage assessment}
\label{sec:damage-assessment}

Assessment is staged because the asset population is large and most of it is
unaffected. The first stage screens on maximum settlement and slope, typically
excluding buildings predicted to settle less than roughly 10\,mm and to tilt
less steeply than 1 in 500.
The second stage assesses buildings that pass that screen against greenfield
movements, treating the building as though it conformed to the ground
movements with no stiffness of its own. The third stage examines the remainder
in detail, accounting for the building's own stiffness, its foundation type,
its condition and its orientation to the trough.

The staging is not merely an economy; it works because of the shape of the
problem. Figure~\ref{fig:settlement-screening} shows that shape along an
illustrative corridor: settlement above 5\,mm is confined to a band within
about 15\,m of the centreline, and only 49 buildings in a population of three
thousand, fewer than two in a hundred, exceed the 10\,mm settlement criterion
of the first stage. The bands
in that figure are settlement magnitudes and not damage categories. The screen
is deliberately crude; the assessment that follows it is not.

\begin{figure}[htbp]
\centering
\datafigure[0.94\textwidth]{ch26/damage-map}
\par\smallskip
\datafigure[0.68\textwidth]{ch26/damage-categories}
\caption[First-stage settlement screening across a building population]{%
What the first stage screens. Above: predicted settlement band by position for
three thousand buildings along an illustrative twin-tunnel corridor, plotted by
chainage and by offset from the centreline. Below: the same population by band.
The bands are magnitudes of predicted settlement and not damage categories ---
they carry no building stiffness, no foundation type and no orientation to the
trough --- and their only purpose is to decide which buildings are assessed
further. Synthetic demonstrator data (Appendix~\ref{app:software}).}
\label{fig:settlement-screening}
\end{figure}

The damage criterion itself is limiting tensile strain. The idealisation of
\citet{burland1974settlement} treats the building as a deep beam undergoing both bending and shear.
The building is divided at the point of inflection of the trough into a
hogging length and a sagging length, each is assessed separately, and the
quantity each consumes is its deflection ratio --- the maximum relative
deflection divided by that length. From it a maximum tensile strain is derived
and compared against thresholds
corresponding to damage categories --- negligible, very slight, slight,
moderate, severe, very severe. The threshold strains are given in
Appendix~\ref{app:constants}, and the full sequence in
Section~\ref{sec:seq-settlement}. The categories are defined by the ease of
repair, which is what an owner cares about, rather than by structural capacity.

\citet{boscardin1989building} extended this by combining the distortion of the
building with the
horizontal strain from the ground movement field, producing an interaction
chart with angular distortion on the horizontal axis, horizontal strain on the
vertical, and damage category contours across it. The extension matters
because horizontal strain shifts many
buildings up a damage category, and this is the step most often omitted in
practice; the horizontal strain itself follows from the expression for
$\varepsilon_h$ in Section~\ref{sec:movement-field}.

Building stiffness is where the greenfield assumption is conservative, for the
deflection ratio and horizontal strain specifically. A real
building does not conform to the ground; it resists, modifying the trough
beneath itself and experiencing less deflection than the greenfield curve
implies, at the price of carrying in its own structure the load that resistance
requires. Relative stiffness modification factors~\citep{potts1997structures}
quantify the reduction, and it
can be substantial for a stiff building on a shallow trough. The factors were
derived from plane-strain analyses of an elastic beam on the ground surface,
with relative bending and axial stiffnesses as the inputs, and they are no
better than the stiffness assumed. Applying them requires care about openings, since a facade with a high window
ratio is far less stiff than its dimensions suggest, and about the building's
orientation --- a building at an angle to the tunnel is subject to a more
complex deformation than a two-dimensional assessment represents. They are
unsafe where the stiffness is assumed rather than demonstrated, as it often is
for the buildings reaching the third stage, and for weak or poorly
tied masonry, which redistributes movement through its joints rather than
resisting it; such a building is assessed on greenfield movements.

\subsection{Foundations and buried assets}
\label{sec:foundations-buried}

Foundation type changes the response. Strip and raft foundations follow the
ground closely and are assessed by the methods above. Piled foundations are
governed by the settlement at the pile toe rather than at the surface, which
requires the subsurface movement field. A pile whose toe is above the tunnel
axis may settle more than the surface; a pile founded below the zone of
influence may not settle at all, in which case the building experiences
differential settlement between piled and non-piled elements. Tunnelling
beside a pile also reduces its shaft friction, transferring load to the toe.

Buried services are assessed on joint rotation for jointed pipelines and on
strain for continuous ones. The governing uncertainty is rarely the analysis:
it is the condition and position of the service. Records of buried
infrastructure in older cities are incomplete and frequently wrong, and a
cast-iron main of unknown age and condition is not amenable to a strain
calculation with any real confidence. Survey, condition assessment and, where
critical, replacement in advance are more reliable than analysis.

Operating railways impose the tightest limits of all. Track twist, cant and
alignment tolerances are measured in millimetres over defined base lengths,
and exceeding them requires a possession to correct. The limits are
maintenance limits rather than safety limits, and they routinely govern the
design of a tunnel passing beneath a railway.

\section{Assessing and protecting assets}
\label{sec:assessing-protecting}

\subsection{Existing tunnels and excavation interaction}
\label{sec:existing-tunnels-interaction}

Tunnelling near an existing tunnel loads a structure whose condition and
capacity are usually poorly known. The new excavation causes the existing
lining to ovalise, its joints to open or close, and its alignment to move. A
segmental lining tolerates ovalisation better than a cast-iron or masonry one,
and a masonry tunnel from the nineteenth century may have little reserve at
all.

Proximity criteria based on clear separation in diameters are a useful screen
and no more. What governs is the movement imposed at the existing structure
and its capacity to accommodate it, so the assessment needs the subsurface
movement field and a structural assessment of the existing tunnel --- often
the harder of the two, because the as-built details are unknown.

Twin tunnels interact with each other. The second drive imposes additional
movement, and it does so on ground already disturbed and loosened by the
first, so its volume loss is typically higher for the same construction
control. Simple superposition of two greenfield troughs therefore
underestimates the total, and modification factors are applied to the second
drive. Pillar width and construction sequence both matter, and a design that
assumes symmetric behaviour from two nominally identical drives has ignored
the mechanism.

\subsection{Subsidence over voids and legacy workings}
\label{sec:void-subsidence}

A different mechanism operates above voids. Where a cavity exists --- an
abandoned mine working, a karst cavern, a poorly backfilled shaft --- the
ground above it can migrate downward progressively, the void working upward as
material falls into it. Because collapsed material bulks, the migration
arrests once the bulked volume fills the void. Collapse of a thickness $h$ of
roof bulks to fill a void of height $H_v$ when $(B-1)\,h = H_v$, where $B$ is
the bulking factor, so that $h = H_v/(B-1)$. The result is acutely sensitive
to a parameter that is itself poorly known: a 3\,m void arrests after about
10\,m of migration at $B = 1.30$ and after about 20\,m at $B = 1.15$. If the arrest height
is less than the cover, nothing appears at the surface; if it exceeds the
cover, a sinkhole appears, often suddenly and long after the working was
abandoned.

The crown pillar between a working and the surface is therefore the quantity
of interest, and its adequacy depends on the span of the void, the strength
and bedding of the overlying rock, and time --- deterioration of a crown
pillar continues for decades.

Investigating legacy voids begins with records research, which is frequently
the most productive step and often the only one that identifies workings at
all. Probe drilling on a grid and geophysics follow, and both have low
detection reliability for small voids. Treatment is by grout filling, by
capping, or by structurally bridging the affected area, and residual risk
remains after all of them, because the investigation cannot demonstrate that
every void has been found.

\subsection{Protection and monitoring}
\label{sec:protection-monitoring}

Where predicted damage is unacceptable, the responses fall in a clear order of
preference. Reducing the ground movement at source --- better face pressure
control, improved grouting, a different alignment or depth --- is the most
effective and the cheapest, and it is frequently dismissed too early because
the alignment is treated as fixed. Ground treatment between the tunnel and the
structure (Chapter~\ref{ch:ground-improvement}) intercepts the movement.
Compensation grouting actively replaces the lost volume in response to
monitoring. Structural strengthening or underpinning of the building accepts
the movement and makes the structure able to tolerate it.

Compensation grouting deserves its reputation and its caveats. It is effective
and it is expensive. It requires access between the tunnel and the structure,
real-time monitoring and a decision process able to act within hours, and it
cannot protect against a sudden loss. It should be planned as a system,
with the sleeve pipes installed and the response established before tunnelling
approaches, not mobilised after settlement has appeared.

Monitoring of third-party assets closes the loop, and its design is as much
contractual as technical. Instruments on other people's property require
agreements; trigger levels become commitments; and exceedance triggers
obligations. The condition survey carried out before work begins is the single
most valuable document in the process, because without it every pre-existing
crack becomes a claim. It should be thorough, photographic, dated and agreed
with the owner --- and it costs a fraction of what the disputes it prevents would cost.
\section{Instrumentation and observational design}
\label{sec:instrumentation-observational}

\subsection{Purpose and the observational method}
\label{sec:purpose-observational}

Monitoring serves five purposes, and being explicit about which applies is
what determines the design of the array. It \emph{verifies design
assumptions}, confirming that the ground is behaving as predicted. It
\emph{controls construction}, giving the face team the information to adjust
support or sequence. It \emph{protects third parties}, demonstrating that
movement limits on adjacent assets are being respected. It provides a
\emph{contractual record} of conditions and behaviour. And it builds
\emph{knowledge} for future projects.

The observational method makes the first three into a design philosophy rather
than a check. The method was set out by \citet{peck1969observational} and
later codified in Eurocode~7, and its requirements can be compressed to four
preconditions: the range of possible behaviour must be
assessed and the acceptable limits established \emph{before} construction; a
plan of monitoring capable of detecting departure from the predicted behaviour
must be devised; the response time of the instruments and the analysis must be
short compared with the rate at which the situation could develop; and
contingency measures must be prepared and capable of being implemented in the
time available.

Every one of those is routinely violated, and the fourth most often. A
monitoring regime with trigger levels but no pre-agreed, resourced contingency
is not the observational method; it is an alarm with nothing behind it. The
method can be applied from the outset, with the design based on the most
probable conditions and contingency held for the unfavourable case, or as a
recovery when behaviour departs from expectation --- but the recovery
application is much weaker, because the preparation was not done.

\subsection{Instrument types}
\label{sec:instrument-types}

\emph{Deformation} instruments are the backbone. Convergence arrays measure
the change in chord lengths across the excavation and are cheap, robust and
directly meaningful. Borehole extensometers measure displacement at several
depths into the rock, which distinguishes movement of the whole rock mass from
loosening of the immediate skin --- a distinction convergence alone cannot
make. Inclinometers measure lateral profile in a borehole and are the standard
instrument for retaining walls. Settlement cells and magnetic extensometers
measure vertical movement at depth.

\emph{Load and stress} instruments are less reliable and should be specified
with that in mind. Load cells on bolts and anchors measure a well-defined
quantity and work well. Pressure cells embedded in shotcrete or between lining
and ground attempt to measure stress in a medium whose stiffness they alter by
being present, and their absolute readings should be treated with caution;
their trends are more trustworthy than their values. Strain gauges in a lining
measure strain reliably and require an assumed modulus to become stress --- for
young shotcrete, that modulus is changing as the reading is taken.

\emph{Groundwater} instruments complete the set. Standpipe piezometers are
simple and slow; vibrating wire piezometers respond quickly and can be
installed in multiple horizons in one borehole. Response time depends on the
permeability of the ground, and in low-permeability clay a standpipe may lag
the true pressure by weeks --- long enough to give a comfortable reading during
exactly the period when the pressure is changing.

\subsection{Surface and remote monitoring}
\label{sec:surface-remote}

Third-party protection is monitored from the surface. Precise levelling
remains the reference method; automated total stations tracking prism arrays
provide continuous measurement without access; tiltmeters detect rotation of
structures; crackmeters track existing cracks whose behaviour is often the
most legible indicator a building offers.

What is achievable in a live urban environment is less than laboratory figures
suggest. Automated total stations are affected by atmospheric refraction,
temperature gradients and instrument stability, and sub-millimetre precision
is claimed more often than it is delivered. Since trigger levels are
frequently set at a few millimetres, the measurement precision and the trigger
threshold can be uncomfortably close, and the defence is redundancy and
trend analysis rather than reliance on any single reading.

Remote sensing extends coverage. Satellite radar interferometry measures
wide-area settlement over time without any site installation and is
particularly effective for drawdown settlement, which is too broad for a
conventional array. Terrestrial laser scanning and photogrammetry measure
profile and convergence over whole surfaces rather than at points, and are the
practical route to measuring overbreak and shotcrete thickness. Ground
penetrating radar locates voids behind a lining, and computer vision detects
and classifies lining defects from imagery.

Wide coverage is not the same as a wide-area answer. A radar record is a rate
measured against a reference frame over an archive period, and the movement the
works caused is the difference between what the record shows during
construction and what the ground was already doing
(Figure~\ref{fig:remote-baseline}). Where no pre-construction archive has been
processed that difference cannot be formed, and the whole of the observed
movement is attributed to the excavation by default.

\begin{figure}[htbp]
\centering
\datafigure[0.94\textwidth]{ch26/insar-baseline}
\par\smallskip
\datafigure[0.94\textwidth]{ch26/settlement-attributable}
\caption[Wide-area radar: the baseline and what is attributed to the works]{%
Wide-area coverage and the interpretation it requires. Above: line-of-sight
rates over an illustrative corridor before excavation begins, which is the
movement field the works inherit. Below: the same buildings once that baseline
is removed, which is the movement for which the works are answerable. The lower
panel is derived rather than measured, and it remains conditional on the archive
period the baseline was fitted over, on the reference frame and on the passage
window assumed for each drive. Synthetic demonstrator data.}
\label{fig:remote-baseline}
\end{figure}

Each of these is a data stream into the model of
Chapter~\ref{ch:digital-twin} rather than a standalone record, and their value
lies in updating the model that the design decisions rest on.

\section{Interpretation and design updating}
\label{sec:interpretation-updating}

\subsection{Data acquisition and array design}
\label{sec:acquisition-array}

Automated acquisition, telemetry and dashboards have removed the delay between
measurement and availability, which was historically the limiting factor in
applying the observational method. What they have introduced is data volume:
a project can now generate more readings than anyone examines, and the
discipline has shifted from collecting data to filtering it.

Array design should start from the design question rather than from the
instrument catalogue. For each assumption the design depends on, the questions
are what would indicate that it is wrong, what instrument would detect that,
where it must be installed and how often it must be read. Instruments that
survive this test earn their place; those that do not are recording for its
own sake.

Three practical requirements follow. Monitoring sections should be spaced to
capture variation along the drive, with additional sections at every feature
of interest --- fault crossings, low cover, third-party assets, changes in
ground class. Instruments must be installed close enough to the face to
capture the deformation that matters, which for convergence means within a
fraction of a diameter, because deformation missed before installation is lost
permanently. And baselines must be established before construction begins; a
monitoring programme that starts after the works have started has no datum,
and its readings are uninterpretable.

\subsection{Trigger levels and interpretation}
\label{sec:triggers-interpretation}

Trigger levels convert measurements into decisions, and the conventional
green--amber--red scheme works when the thresholds derive from the design
prediction rather than from convenience. The amber level should correspond to
behaviour departing from prediction while remaining acceptable, prompting
increased monitoring and review; the red level to behaviour approaching the
limit of what the design can tolerate, prompting the contingency measure.
Setting them by rounding down a calculated capacity produces triggers that
either fire constantly or never fire at all.

Each level needs an owner, an escalation path and an action. Who is informed,
who decides, within what time, and what happens --- specified in advance, not
determined during the event.

Interpretation is where experience is required. Instrument drift, temperature
effects, construction disturbance and simple malfunction all produce readings
that resemble ground movement. Cross-checking between instruments of different
type is the reliable defence: convergence and extensometer readings that
disagree indicate an instrument problem. Agreement removes independent
malfunction but not a common cause, since temperature and movement of the
datum act on both alike, so confidence in a trend calls for a third,
physically independent observation.

The quantity to watch is usually the rate rather than the magnitude.
Displacement that grows and decelerates is a system finding equilibrium;
displacement that grows at constant or increasing rate is a system that has
not. A tunnel can accommodate large convergence safely if it is converging
towards a limit, and a much smaller convergence at a steady rate is the more
serious observation.

\subsection{Design updating and common failures}
\label{sec:updating-failures}

Monitoring closes the design loop by feeding back into the ground model and
the support design. Measured behaviour calibrates the numerical model
(Chapter~\ref{ch:numerical-modelling}), and the calibrated model predicts the
remainder of the drive far better than the original. Ground class assignments
ahead of the face can be revised, support classes adjusted, and the prediction
against which the next section is compared updated.

Back-analysis carries the non-uniqueness caveat already noted: several
parameter combinations reproduce one measured quantity, so calibration should
use several --- convergence, extensometer profile and support load together.

The recurring failures are worth stating plainly, because they are failures of
process rather than of instrumentation. Instruments installed too late to
capture the initial response. No pre-construction baseline. Data collected and
never examined. Trigger levels with no owner and no contingency. Monitoring
specified as a deliverable to be provided rather than as a decision tool to be
used. And the design team disbanded before the evidence arrives, so that
nobody with the knowledge to interpret the readings is still on the project
when they are taken.

The last is the most consequential and the easiest to prevent: the
observational method requires the designer to remain engaged through
construction, and that engagement has to be in the contract rather than in
anyone's goodwill.

\section{Remote sensing and as-built geometry}
\label{sec:remote-sensing-geometry}

The instruments of Section~\ref{sec:instrument-types} measure points. Laser scanning and
photogrammetry measure surfaces, returning a dense record of the excavation as
it actually is rather than at the small number of locations an array happens
to occupy. That difference in kind changes what can be asked: not only how
much a section has converged, but where the profile departs from the design,
how much material was taken beyond it, and whether the lining is as thick as
specified.

The technique is now routine, and the data is frequently captured and then
used for nothing beyond a visual record. The reason is usually that nobody
decided in advance which engineering question the survey was to answer.

\subsection{Acquisition and registration}
\label{sec:point-cloud-acquisition}

Two methods dominate. Terrestrial laser scanning measures range directly and
performs well in the dark, wet, dusty conditions of a heading; it is fast, and
accuracy at tunnel scale is typically a few millimetres. Photogrammetry
recovers geometry from overlapping images and needs lighting and texture, but
the equipment is cheap, it carries colour information useful for face mapping
and defect recording, and it can be flown or hand-held. Mobile systems mounted
on a vehicle or on the machine trade some accuracy for the ability to survey
without stopping the works, which in an operating drive is often the decisive
consideration.

Registration is where the accuracy is won or lost. A scan is measured in the
instrument's own frame, and converting it to project coordinates requires
control --- surveyed targets, or overlap with a previous scan, or both. The
resulting registration error is systematic rather than random, it accumulates
along a drive when scans are chained to one another rather than tied to
control, and it is invisible in the point cloud itself, which looks equally
convincing whether correctly registered or not. Any statement of accuracy that
quotes the instrument's ranging precision without the registration error is
describing the wrong quantity.

The adequacy test is against the question. Detecting overbreak of a few
hundred millimetres tolerates registration error of ten; detecting convergence
of two millimetres between epochs does not, and that case demands the
discipline of Section~\ref{sec:change-detection}.

\subsection{As-built against as-designed}
\label{sec:as-built}

Comparing the surveyed surface with the design profile produces several
quantities at once, which is the technique's principal economy.

Overbreak and underbreak are measured directly, over the whole surface rather
than at the section a profile gauge happens to check. Overbreak is a cost and
a carbon quantity through the volume of shotcrete or concrete needed to fill
it, and a quality indicator for the perimeter blasting practice of
Chapter~\ref{ch:construction}: a systematic pattern points at the drilling and
charging rather than at the ground. Underbreak is an immediate construction
problem, since the clearance envelope must be achieved.

Applied sprayed concrete thickness follows from differencing the scan before
and the scan after spraying, the verification route introduced with sprayed
concrete in Chapter~\ref{ch:support}. This is worth emphasising because thickness is
otherwise verified by depth pins at a handful of points, and the difference
between specified and achieved thickness is a recurring durability problem
(Chapter~\ref{ch:linings}). The same differencing gives the volume placed,
which reconciles against delivery records and exposes rebound losses that are
otherwise estimated.

Finally, the clearance envelope can be checked against the surveyed surface
rather than against the design, which is the check that actually matters
before fit-out. Where the structure is handed over with this record, the
operator inherits a dimensioned account of what was built, which is the
starting point for the inspection regime of Chapter~\ref{ch:assets}.

\subsection{Change detection over time}
\label{sec:change-detection}

Repeat surveys can measure movement, but the analysis is more demanding than
the single-epoch case, and treating one as the other produces confident
nonsense.

The difficulty is that a point in one cloud has no counterpart in the next:
the surfaces are sampled independently, so differencing must compare surfaces
rather than points, and the comparison is contaminated by registration error
between epochs. That error behaves like a rigid movement of the whole scan,
which is precisely the signature that a genuine uniform convergence would
also produce. Separating them requires stable reference surfaces inside the
scanned volume --- a completed lining section, a rock face outside the zone of
influence --- against which each epoch is registered.

Within that constraint the method is powerful. It detects movement anywhere on
the surface rather than at instrumented points, which is what allows an
unexpected mechanism to be seen at all; a conventional array measures only
where somebody anticipated a problem. It suits slow processes with no fixed
location, such as progressive deterioration or subsidence over legacy
workings. And it produces a record that can be revisited when a question
arises years later, which point measurements cannot.

It does not replace instrumentation. Scanning is episodic where an array is
continuous, and the observational method of
Section~\ref{sec:purpose-observational} depends on a measurement
frequency high enough to act on before the ground does. The sound arrangement
uses both: continuous instrumentation at the sections identified as critical,
periodic scanning to establish that nothing is happening where nobody is
looking.

\chapterreferences
